\documentclass{article}

% =========================================================
% (1) CORE PACKAGES
% =========================================================
% bifurcationatlas:
%   Main package providing the atlas system.
%
%   Core user-level commands:
%   - \initAtlas        : define reusable atlas templates
%   - \setAtlas         : select active atlas template
%   - \baNewRegion      : define labeled regions with metadata
%   - \baAtlasFigure    : render central bifurcation diagram
%   - \baRadialAtlas    : internal radial layout engine
%   - \baSetup          : global configuration interface
%
% graphicx:
%   Standard image inclusion (used for center diagrams)
%
% geometry:
%   Page layout control (ensures symmetric margins for figures)
% =========================================================

\usepackage{bifurcationatlas}
\usepackage{graphicx}
\usepackage[margin=2.5cm,centering]{geometry}

\graphicspath{{./fig/}}

% =========================================================
% (2) GLOBAL COLOR DEFINITIONS
% =========================================================
% Semantic color system:
%
% Colors encode dynamical meaning rather than aesthetics:
% - Stability classes: Stable / Unstable / Saddle / Focus
% - Equilibrium families: HHEqui, HLEqui, MEqui, LLEqui, TEqui
%
% This allows model-independent visualization where
% color = semantic state, not arbitrary styling.
% =========================================================

% --- stability palette ---
\definecolor{Stable}{RGB}{0,110,0}
\definecolor{Unstable}{RGB}{255,77,77}
\definecolor{Saddle}{RGB}{255,165,0}
\definecolor{Focus}{RGB}{102,102,255}
\definecolor{Periodic}{RGB}{26,26,255}

% --- equilibrium-type palette ---
\definecolor{HHEqui}{RGB}{0,114,178}
\definecolor{HLEqui}{RGB}{230,159,0}
\definecolor{MEqui}{RGB}{0,158,115}
\definecolor{LLEqui}{RGB}{204,121,167}
\definecolor{TEqui}{RGB}{213,94,0}

% =========================================================
% (3) ATLAS TEMPLATE SYSTEM (GEOMETRIC DEFINITIONS)
% =========================================================
%
% Each \initAtlas defines a reusable geometric layout:
% - region identifiers (e.g., LL, HL, M, HH, C)
% - TikZ drawing instructions per region
% - semantic mapping to colors via \baColor{<key>}
%
% This separates:
%   geometry (structure)
%   from
%   data (region labeling)
% =========================================================

% ---------------------------------------------------------
% TEMPLATE: Cartel (5-region structured system)
% ---------------------------------------------------------
\initAtlas{Cartel}{LL,HL,M,HH,C}{
  % base blocks (semantic color injection)
  \fill[\baColor{LL}] (0,0) rectangle (1,1);
  \fill[\baColor{HL}] (1,0) rectangle (2,1);
  \fill[\baColor{M}]  (0,1) rectangle (1,2);
  \fill[\baColor{HH}] (1,1) rectangle (2,2);

  % refinement region (high-resolution structure)
  \fill[\baColor{C}] (0,0) rectangle (0.5,0.5);

  % structural grid overlay
  \draw[line width=2pt,black] (0,2)--(0,0)--(2,0);
  \draw (0,1)--(2,1);
  \draw (1,0)--(1,2);

  % refinement grid
  \draw (0,0.5)--(0.5,0.5);
  \draw (0.5,0)--(0.5,0.5);

  % highlight boundary
  \draw[line width=2pt,red] (0,0.04)--(0,2);
}

% ---------------------------------------------------------
% TEMPLATE: SimpleState (2-region system)
% ---------------------------------------------------------
\initAtlas{SimpleState}{SL,SR}{
  \fill[\baColor{SL}] (0,0) rectangle (1,2);
  \fill[\baColor{SR}] (1,0) rectangle (2,2);
}

% =========================================================
% (4) GLOBAL CONFIGURATION
% =========================================================
%
% These parameters define global rendering behavior:
% - region scale  : scaling of individual region diagrams
% - region radius : radial placement radius in atlas layout
% - figure scale  : scaling of central bifurcation diagram
% - debug         : enables overlay diagnostics
% =========================================================

\baSetup{
  region scale = 1.25,
  region radius = 6.75,
  figure scale = 0.35,
  debug = false,
}

% =========================================================
% (5) DOCUMENT START
% =========================================================
\begin{document}

% =========================================================
% (6) ATLAS EXAMPLE 1: SimpleState (arabic labels)
% =========================================================
\begin{bifurcationatlas}[region scale = 1.5]
\label{ba:simplestate}
\setAtlas{SimpleState}
\centering

% Define regions (data layer), default color is transparent
\baNewRegion[caption={Region 1},label={ba:R1}]{R1}{SL=Unstable}
\baNewRegion[caption={Region 2},label={ba:R2}]{R2}{SL=Stable,SR=Unstable}
\baNewRegion[caption={Region 3},label={ba:R3}]{R3}{SL=Stable,SR=Focus}

% Central bifurcation diagram + radial overlay
\baAtlasFigure[caption={Cartel Diagram},label=ba:center]{Center_Corridor_Region}{R1,R2,R3}

\end{bifurcationatlas}

% =========================================================
% (7) ATLAS EXAMPLE 2: SimpleState (roman labels)
% =========================================================
\begin{bifurcationatlas}[caption label=roman]
\label{ba:simplestate2}
\setAtlas{SimpleState}
\centering

\baNewRegion[caption={Region A},label={ba2:R1}]{R1}{SL=Stable,SR=Focus}
\baNewRegion[caption={Region B},label={ba2:R2}]{R2}{SL=Unstable,SR=Unstable}
\baNewRegion[caption={Region C},label={ba2:R3}]{R3}{SL=Periodic,SR=Focus}

\baAtlasFigure[caption={Cartel Diagram Second},label=ba2:center]{Center_Corridor_Region}{R1,R2,R3}

\end{bifurcationatlas}

% =========================================================
% (8) ATLAS EXAMPLE 3: Cartel (upper roman labels)
% =========================================================
\begin{bifurcationatlas}[caption label=Roman]
\label{ba:cartel}
\setAtlas{Cartel}
\centering

% Dense semantic mapping example (multi-parameter system)
\baNewRegion[caption={Cartel Region 1},label={ba:cartel_R1}]
  {R1}{LL=LLEqui,HL=HLEqui,M=HHEqui,HH=HHEqui,C=TEqui}

\baNewRegion[caption={Cartel Region 2},label={ba:cartel_R2}]
  {R2}{LL=LLEqui,HL=HLEqui,M=MEqui,HH=HHEqui,C=TEqui}

\baNewRegion[caption={Cartel Region 3},label={ba:cartel_R3}]
  {R3}{LL=LLEqui,HL=HLEqui,M=MEqui,HH=MEqui,C=TEqui}

\baNewRegion[caption={Cartel Region 4},label={ba:cartel_R4}]
  {R4}{LL=LLEqui,HL=HLEqui,M=HLEqui,HH=HLEqui,C=HLEqui}

\baNewRegion[caption={Cartel Region 5},label={ba:cartel_R5}]
  {R5}{LL=LLEqui,HL=LLEqui,M=LLEqui,HH=TEqui,C=TEqui}

\baNewRegion[caption={Cartel Region 6},label={ba:cartel_R6}]
  {R6}{LL=LLEqui,HL=LLEqui,M=TEqui,HH=TEqui,C=HLEqui}

% Central diagram for full atlas
\baAtlasFigure[caption={Cartel Bifurcation Diagram (roman)},label={ba:cartel_center}]{bifurcation}{R1,R2,R3,R4,R5,R6}

% Global caption for interpretation
\caption{Name-based region test using semantic color mappings.}

\end{bifurcationatlas}

% =========================================================
% (9) ILLUSTRATIVE DISCUSSION SECTION (FAKE EXAMPLE TEXT)
% =========================================================
\section{Illustrative Comparison of Atlas Constructions}

To demonstrate the flexibility of the bifurcation atlas framework, we compare several structurally distinct atlas configurations introduced in this example.

\ref{ba:simplestate2} shows a simple two-region decomposition using the \texttt{SimpleState} template. The radial arrangement in Subfigure~\ref{ba2:center} illustrates how even minimal structural complexity already yields a meaningful separation of dynamical regimes.

In contrast, \ref{ba:cartel} presents a higher-dimensional semantic decomposition using the \texttt{Cartel} template. Here, multiple equilibrium classes interact across six regions with heterogeneous stability properties. A representative region-level structure can be seen in Subfigure~\ref{ba:cartel_R2}, which highlights the coexistence of unstable and periodic dynamics within a single parameter block.

The radial organization further emphasizes global coherence: individual regions such as \subref{ba:cartel_R5} are not isolated visual elements but contribute to a consistent global phase organization around the central bifurcation diagram.

Overall, the comparison between the SimpleState and Cartel atlases illustrates the scalability of the framework from minimal binary decompositions to rich multi-class semantic partitions. This is particularly visible when comparing the compact structure of \ref{ba:simplestate2} with the densely structured atlas in \ref{ba:cartel}.
% =========================================================
% END DOCUMENT
% =========================================================

\end{document}